\subsection{Learning Without Rational Expectations}
\label{section:fc_paradigms}

This chapter studies world models in reinforcement learning, learned forecasters of state transitions and rewards used by an agent to plan its actions. In the language of economics a world model is a perceived law of motion that the agent estimates from data and then plans against, an object that economics has studied since at least the 1980s under names such as adaptive learning, fictitious play, and Berk-Nash equilibrium. The subsection that follows briefly catalogues these analogues before the chapter develops the reinforcement-learning line in detail.

\subsubsection{A family of non-rational-expectations learning approaches}
\label{section:fc_paradigm_family}

The first approach is belief-based learning from game theory, in which an agent maintains a representation of opponents' play or of the environment and refines it as new observations arrive. Fictitious play \citep{Brown1951}, its smoothed variant \citep{FudenbergKreps1993}, and stochastic fictitious play under global convergence conditions \citep{HofbauerSandholm2002} all share with model-based reinforcement learning the basic move of starting from a possibly wrong belief and updating it from data, and they ask the same kinds of convergence question (whether beliefs concentrate, whether they diverge, whether they select among multiple equilibria). The textbook treatments in \citet{FudenbergLevine1998} and \citet{Young2004} carry the main results, and the parallels collected below are in that same family rather than formal equivalences.

The second is recursive least squares adaptive learning, developed by \citet{MarcetSargent1989} and synthesized in \citet{EvansHonkapohja2001}. In this approach the agent holds a perceived linear law of motion and updates its coefficients from realized data; convergence to a rational-expectations equilibrium holds when an associated ordinary differential equation is locally stable at the fixed point and fails otherwise. The construction is single-rule and on-policy, and it contains no explicit exploration mechanism, which sharply constrains the regions in which it converges.

The third is genetic algorithm learning \citep{Arifovic1994, Arifovic1995}. A population of binary-encoded decision rules evolves under fitness-proportional selection, crossover, mutation, and an election operator that admits offspring only when they would have outperformed their parents on the previous environment. The augmented genetic algorithm converges to the rational-expectations equilibrium of the cobweb in both stable and unstable parameter regions, including those in which recursive least squares diverges, with exploration provided by population diversity rather than by a noise term on a single policy.

The fourth is evolutionary heuristic switching \citep{BrockHommes1997}, in which agents choose among a finite set of forecasting heuristics by a logit on past performance with intensity-of-choice parameter $\beta$. As $\beta$ rises from zero the joint dynamics trace a period-doubling cascade from a stable steady state to high-dimensional attractors, and the source of complexity is the discrete selection rule itself rather than the model class or any exploration operator.

The fifth is misspecified Bayesian learning, formalized as Berk-Nash equilibrium by \citet{EspondaPouzo2016}. The agent's subjective model class need not contain the truth; beliefs concentrate on parameters minimizing Kullback-Leibler divergence to the data-generating process, and at equilibrium the strategy is optimal given the belief while the belief is best-fit given the strategy. Misspecification is permanent rather than a transient phase as in the four approaches above.

The parallel collected here is architectural rather than an equivalence claim; adaptive learning asks whether beliefs converge while model-based reinforcement learning asks whether a learned transition model improves control, and the success criterion that follows the two literatures is correspondingly different.

\label{section:fc_convergence_question}%
The five approaches share a common formal vocabulary, drawn from stochastic approximation and fixed-point analysis, and they address related but distinct convergence questions. Single-agent model-based reinforcement learning, the line this chapter touches on next, asks whether a policy approaches optimality on the true environment as the learned model and the policy are updated together. Multi-agent learning in agent-based models, including \citet{MarimonMcGrattanSargent1990}, \citet{KrusellSmith1998}, and \citet{Calvano2020}, asks which rational-expectations equilibrium the joint dynamics select when several exist, and whether the selected equilibrium matches the lab-experimental or empirical attractor. Both questions concern the limiting behavior of non-rational-expectations learning dynamics, the machinery is similar, and the objects differ; one is the optimality gap of a single policy, the other is the basin of attraction of a system of interacting learners. The machinery for both traces back to the algorithmic background developed in \S\ref{section:rl_algorithms}. With this brief catalogue in place, the chapter turns next to the two 1990 origin papers that opened the model-based line, one from the reinforcement-learning tradition and one from the recurrent-neural-network tradition.
